\section{本讲小结}

电流是有符号的净电荷通量，电流密度须同时考虑载流子的电荷和漂移方向。稳恒状态保证无分支通道各截面电流相同，截面与材料性质的变化仍可使电流密度和电场变化。载流导体需要内部电场，不能套用静电平衡的零场条件。

整体电阻来自局域欧姆规律与实际几何的结合。相同电流依次穿过的薄层按串联积分，共用两等势端面的通道按电导相加。径向圆柱、完整球面和大地半空间对应不同的有效面积，求场与电势时应先明确对称性和边界，再选择微元。局部场乘步长只在场变化足够小时近似成立。

漂移与碰撞模型描述载流子的平均输运。弛豫时间、固定碰撞间隔和电路传播时间分别表征不同过程；制动后的渐近速度、有限位移以及惯性电荷分离也须按相应受力条件求解。若电荷积累建立了不可忽略的附加场，原先忽略电场的模型就应停止使用。

电动势表征非静电作用提供的单位电荷能量，端电压则由电源的充放电方向和内阻共同决定。焦耳热、其他能量转换与储能须逐项列出，最大外功率也不等同于最大效率。温度影响电阻时，应联立电路与热量方程，并检查热平衡的存在条件和稳定性。
